\documentclass[12pt,a4paper]{article}
\usepackage[utf8]{inputenc}
\usepackage[T1]{fontenc}
\usepackage{geometry}
\geometry{left=2.4cm,right=2.4cm,top=2.4cm,bottom=2.4cm}
\usepackage{amsmath,amssymb,amsthm,mathtools}
\usepackage{booktabs}
\usepackage{setspace}
\usepackage[hidelinks]{hyperref}
\usepackage[
  backend=biber,
  style=authoryear,
  natbib=true,
  sorting=nyt
]{biblatex}
\addbibresource{reference.bib}
\onehalfspacing

\newtheorem{definition}{Definition}
\newtheorem{proposition}{Proposition}
\newtheorem{remark}{Remark}

\newcommand{\R}{\mathrm{R}}
\newcommand{\F}{\mathrm{F}}
\newcommand{\E}{\mathbb E}
\newcommand{\1}{\mathbb 1}

\title{Spatial Equilibrium and Optimal Transportation Investment\\under Exogenous Zoning}
\author{Ke Jiang}
\date{July 2026}

\begin{document}
\maketitle

\section{Purpose and timing}

This version treats municipal zoning as exogenous. Let
\[
\mathbf z
=
\{s_i^{\R},s_i^{\F},\bar f_i^{\R},\bar f_i^{\F}\}_{i\in S},
\]
where $s_i^u$ is the land share assigned to use $u\in\{\R,\F\}$ and $\bar f_i^u$ is the corresponding FAR cap. These policy variables satisfy
\[
s_i^{\R}+s_i^{\F}=1,
\qquad
s_i^u\geq 0,
\qquad
\bar f_i^u>0.
\]
The transportation authority takes $\mathbf z$ as given and chooses network investment. Conditional on zoning and investment, households, firms, developers, rents, wages, OD flows, recursive path choices, and congestion are determined jointly in spatial equilibrium.

\section{Environment and transportation network}

The metropolitan area contains locations $i\in S$ connected by a directed network $(S,E)$. Each location has $\Delta_i>0$ units of land. Exogenous zoning assigns
\begin{align}
\bar L_i^u&\equiv s_i^u\Delta_i,
&
\bar H_i^u&\equiv \bar L_i^u\bar f_i^u,
\qquad u\in\{\R,\F\},
\label{eq:zoning-primitives}
\end{align}
where $\bar L_i^u$ is land zoned for use $u$ and $\bar H_i^u$ is the maximum floor space permitted by the FAR cap.

Investment on edge $e=(k,\ell)$ is $I_{k\ell}\geq0$. Following \citet{Fajgelbaum2020-yv}, edge traffic is $n_{k\ell}$ and travel time is
\begin{align}
\operatorname{time}_{k\ell}
=
\delta_{k\ell}^0+
\delta_{k\ell}^1\frac{n_{k\ell}^{\rho}}{I_{k\ell}^{\beta}},
\qquad
\delta_{k\ell}^1,\rho,\beta>0,
\label{eq:edge-time}
\end{align}
and the multiplicative edge cost is
\begin{align}
d_{k\ell}=\exp\!\left(\kappa\operatorname{time}_{k\ell}\right).
\label{eq:edge-cost}
\end{align}
For any directed network walk $r$ from $i$ to $j$, walk cost is
\begin{align}
\tau_{ijr}
=
\prod_{(k,\ell)\in r}d_{k\ell}.
\label{eq:walk-cost}
\end{align}
Let $A(\mathbf d)$ be the weighted adjacency matrix with
$A_{k\ell}=d_{k\ell}^{-\theta}$ for $(k,\ell)\in E$ and zero otherwise. Assume its spectral radius satisfies $\rho(A)<1$. Then
\begin{align}
G
&\equiv
(\mathbf I-A)^{-1}
=
\mathbf I+A+A^2+\cdots,
\label{eq:recursive-G}\\
\tau_{ij}
&=
G_{ij}^{-1/\theta}.
\label{eq:route-costs}
\end{align}
Following \citet{Allen2022-ms}, this inverse aggregates all recursively generated network walks. Conditional on OD pair $(i,j)$, the probability of selecting walk $r$ and the probability of using edge $(k,\ell)$ are
\begin{align}
\pi_{r\mid ij}
&=
\left(\frac{\tau_{ij}}{\tau_{ijr}}\right)^\theta,
\label{eq:route-share}\\
\pi_{ij}^{k\ell}
&=
\sum_{r\ni(k,\ell)}\pi_{r\mid ij}
=
\frac{G_{ik}A_{k\ell}G_{\ell j}}{G_{ij}}
=
\left(
\frac{\tau_{ij}}
{\tau_{ik}d_{k\ell}\tau_{\ell j}}
\right)^\theta.
\label{eq:edge-use-share}
\end{align}
The final expression includes all recursive walks using the edge; it is not based on a finite, exogenously enumerated route set.

\section{Households}

A household chooses residence $i$, workplace $j$, and a recursively generated network walk $r$. Utility is
\begin{align}
u_{ijr}(\omega)
=
\frac{B_i(N_i,L_i)}{\tau_{ijr}}
\left(\frac{Q_{ij}}{\alpha}\right)^{\alpha}
\left(\frac{h_{ij}^{\R}}{1-\alpha}\right)^{1-\alpha}
v_{ijr}(\omega),
\label{eq:household-utility}
\end{align}
subject to
\[
Q_{ij}+q_i^{\R}h_{ij}^{\R}\leq w_j.
\]
The multiplicative taste shock $v_{ijr}(\omega)$ has a nested
extreme-value structure that generates conditional walk-choice
elasticity $\theta$ and residence--workplace choice elasticity
$\epsilon$.
Local amenities satisfy
\begin{align}
B_i(N_i,L_i)
=
\bar B_i \phi_i(N_i,L_i),
\qquad
\phi_i(N_i,L_i)=N_i^aL_i^b,
\qquad a,b<0.
\label{eq:amenity}
\end{align}
Optimal demands are
\begin{align}
Q_{ij}^*&=\alpha w_j,
&
h_{ij}^{\R*}&=\frac{(1-\alpha)w_j}{q_i^{\R}}.
\label{eq:household-demands}
\end{align}
After aggregating over recursive walks, deterministic indirect utility is
\begin{align}
\widetilde V_{ij}
=
\frac{B_i(N_i,L_i)w_j}
{\tau_{ij}(q_i^{\R})^{1-\alpha}}.
\label{eq:indirect-utility}
\end{align}
The residence--workplace share is
\begin{align}
\pi_{ij}
&=
\frac{T_iE_j\widetilde V_{ij}^{\epsilon}}{\Phi},
&
\Phi
&\equiv
\sum_{m,n\in S}T_mE_n\widetilde V_{mn}^{\epsilon}.
\label{eq:od-share}
\end{align}
Metropolitan ex ante utility is
\begin{align}
\bar U
=
\Gamma\!\left(1-\frac1\epsilon\right)\Phi^{1/\epsilon}.
\label{eq:expected-utility}
\end{align}
Population and employment are
\begin{align}
N_i&=N\sum_j\pi_{ij},
&
L_j&=N\sum_i\pi_{ij}.
\label{eq:population-employment}
\end{align}

\section{Competitive floor-space supply under exogenous zoning}

For use $u\in\{\R,\F\}$, developers use land $L_i^u$ and non-land capital $K_i^u$ to produce
\begin{align}
H_i^u
=
\Xi_i^u(L_i^u)^{\vartheta}(K_i^u)^{1-\vartheta},
\qquad
\Xi_i^u>0,
\quad
\vartheta\in(0,1).
\label{eq:construction}
\end{align}
The price of non-land capital is one. With positive demand and zero outside land use,
\[
L_i^u=\bar L_i^u=s_i^u\Delta_i.
\]
Define actual FAR $f_i^u=H_i^u/L_i^u$. Minimum capital per unit of land needed to produce FAR $f$ is
\begin{align}
c_i^u(f)
=
\left(\frac{f}{\Xi_i^u}\right)^{1/(1-\vartheta)}.
\label{eq:construction-cost}
\end{align}
A competitive developer solves
\begin{align}
\max_{0\leq f\leq\bar f_i^u}
\quad
q_i^u f-c_i^u(f)-r_i^u.
\label{eq:developer-problem}
\end{align}
Let
\begin{align}
\varpi
&\equiv\frac{1-\vartheta}{\vartheta},
&
\mathcal K_i^u
&\equiv
\bar L_i^u(\Xi_i^u)^{1/\vartheta}(1-\vartheta)^{\varpi}.
\label{eq:supply-constants}
\end{align}
The unconstrained FAR and floor-space supply are
\begin{align}
f_i^{u,0}(q_i^u)
&=
(\Xi_i^u)^{1/\vartheta}
\bigl[(1-\vartheta)q_i^u\bigr]^{\varpi},
\label{eq:far-unconstrained}\\
H_i^{u,0}(q_i^u)
&=
\mathcal K_i^u(q_i^u)^{\varpi}.
\label{eq:supply-unconstrained}
\end{align}
Actual supply is therefore
\begin{align}
H_i^u
=
\min\left\{
\bar H_i^u,
\mathcal K_i^u(q_i^u)^{\varpi}
\right\}.
\label{eq:floorspace-supply}
\end{align}
Land rent is the residual
\begin{align}
r_i^u=q_i^uf_i^u-c_i^u(f_i^u).
\label{eq:land-rent}
\end{align}
Zoning is binding when the unconstrained floor-space quantity exceeds $\bar H_i^u$.

\section{Firms and the commercial floor-space block}

A firm at workplace $j$ produces
\begin{align}
Y_j=A_jL_j^{\gamma}(H_j^{\F})^{1-\gamma},
\qquad
A_j=\bar A_jL_j^{\zeta},
\qquad
0<\zeta<1-\gamma.
\label{eq:firm-production}
\end{align}
The agglomeration term is external to the individual firm. Conditional on aggregate employment $L_j$ and commercial floor space $H_j^{\F}$, factor prices are
\begin{align}
w_j
&=
\gamma\bar A_jL_j^{\gamma+\zeta-1}
(H_j^{\F})^{1-\gamma},
\label{eq:wage}\\
q_j^{\F}
&=
(1-\gamma)\bar A_jL_j^{\gamma+\zeta}
(H_j^{\F})^{-\gamma}.
\label{eq:commercial-rent}
\end{align}

The commercial block can be solved analytically conditional on $L_j$. Define
\begin{align}
\mathcal C_j(L_j)
\equiv
(1-\gamma)\bar A_jL_j^{\gamma+\zeta}.
\label{eq:commercial-C}
\end{align}
If the commercial FAR cap is slack, combining \eqref{eq:floorspace-supply} and \eqref{eq:commercial-rent} yields
\begin{align}
H_j^{\F,0}(L_j)
=
\left[
\mathcal K_j^{\F}
\mathcal C_j(L_j)^{\varpi}
\right]^{1/(1+\gamma\varpi)}.
\label{eq:commercial-unconstrained-solution}
\end{align}
Hence the equilibrium commercial quantity and rent conditional on $L_j$ are
\begin{align}
H_j^{\F}
&=
\min\left\{
\bar H_j^{\F},
H_j^{\F,0}(L_j)
\right\},
\label{eq:commercial-quantity-solution}\\
q_j^{\F}
&=
\mathcal C_j(L_j)(H_j^{\F})^{-\gamma},
\label{eq:commercial-price-solution}
\end{align}
and the equilibrium wage follows from \eqref{eq:wage}.

\section{Residential floor-space market}

Aggregate housing expenditure by residents of $i$ is
\begin{align}
\mathcal E_i^{\R}
\equiv
(1-\alpha)N\sum_j\pi_{ij}w_j.
\label{eq:residential-expenditure}
\end{align}
Residential market clearing is
\begin{align}
q_i^{\R}H_i^{\R}=\mathcal E_i^{\R}.
\label{eq:residential-clearing}
\end{align}
The residential block can also be solved analytically conditional on OD shares and wages. If the residential FAR cap is slack,
\begin{align}
q_i^{\R,0}
&=
\left(\frac{\mathcal E_i^{\R}}{\mathcal K_i^{\R}}\right)^{\vartheta},
\label{eq:residential-rent-unconstrained}\\
H_i^{\R,0}
&=
(\mathcal K_i^{\R})^{\vartheta}
(\mathcal E_i^{\R})^{1-\vartheta}.
\label{eq:residential-quantity-unconstrained}
\end{align}
Therefore, conditional on $\{\pi_{ij},w_j\}$,
\begin{align}
H_i^{\R}
&=
\min\left\{
\bar H_i^{\R},
H_i^{\R,0}
\right\},
\label{eq:residential-quantity-solution}\\
q_i^{\R}
&=
\frac{\mathcal E_i^{\R}}{H_i^{\R}}.
\label{eq:residential-price-solution}
\end{align}
When the cap binds, quantity is fixed at $\bar H_i^{\R}$ and rent absorbs the demand response. When it is slack, both quantity and rent respond endogenously.

\section{Traffic and spatial equilibrium}

Given OD shares and conditional recursive-walk choices, edge traffic is
\begin{align}
n_e
=
N\sum_{i,j\in S}\pi_{ij}\pi_{ij}^{e}.
\label{eq:edge-traffic}
\end{align}
The equilibrium is most compactly represented as a fixed point in OD shares and edge traffic. Let
\[
x\equiv(\boldsymbol\pi,\mathbf n),
\]
where $\boldsymbol\pi=\{\pi_{ij}\}_{i,j\in S}$ and $\mathbf n=\{n_e\}_{e\in E}$. For given $(\mathbf I,\mathbf z)$, construct the mapping $\mathcal T(x;\mathbf I,\mathbf z)$ as follows:
\begin{enumerate}
\item Use $\boldsymbol\pi$ to compute $N_i$ and $L_j$ from \eqref{eq:population-employment}.
\item Use $L_j$ and exogenous commercial zoning to compute $H_j^{\F}$, $q_j^{\F}$, and $w_j$ from \eqref{eq:commercial-quantity-solution}--\eqref{eq:wage}.
\item Use $\boldsymbol\pi$ and wages to compute $\mathcal E_i^{\R}$, $H_i^{\R}$, and $q_i^{\R}$ from \eqref{eq:residential-expenditure}--\eqref{eq:residential-price-solution}.
\item Use $(N_i,L_i)$ to compute amenities from \eqref{eq:amenity}.
\item Use $(\mathbf n,\mathbf I)$ to compute edge costs, composite OD travel costs, walk shares, and $\pi_{ij}^e$ from \eqref{eq:edge-time}--\eqref{eq:edge-use-share}.
\item Update OD shares using \eqref{eq:indirect-utility}--\eqref{eq:od-share} and update edge traffic using \eqref{eq:edge-traffic}.
\end{enumerate}

\begin{definition}[Spatial equilibrium conditional on zoning and investment]
For exogenous zoning $\mathbf z$ and investment $\mathbf I$, a spatial equilibrium is a fixed point
\begin{align}
x^*(\mathbf I;\mathbf z)
=
\mathcal T\bigl(x^*(\mathbf I;\mathbf z);\mathbf I,\mathbf z\bigr),
\label{eq:spatial-fixed-point}
\end{align}
together with the prices, wages, floor-space quantities, recursive-walk shares, and amenities recovered from steps 1--5 above.
\end{definition}

Equivalently, write all equilibrium conditions as
\begin{align}
F(x,\mathbf I;\mathbf z)=0.
\label{eq:equilibrium-system}
\end{align}
If $F_x$ is nonsingular at an equilibrium, the local investment response of the equilibrium is
\begin{align}
\frac{\partial x^*}{\partial I_e}
=
-F_x^{-1}F_{I_e}.
\label{eq:equilibrium-derivative}
\end{align}
This derivative contains all general-equilibrium responses through construction, rents, wages, sorting, route choice, and congestion.

\section{Transportation authority}

The transportation authority takes zoning $\mathbf z$ as exogenous and anticipates the conditional spatial equilibrium. Define reduced-form welfare
\begin{align}
\widehat U(\mathbf I;\mathbf z)
\equiv
\bar U\bigl(x^*(\mathbf I;\mathbf z),\mathbf I;\mathbf z\bigr).
\label{eq:reduced-welfare}
\end{align}
The authority solves
\begin{align}
\max_{\mathbf I\geq0}
\quad
\widehat U(\mathbf I;\mathbf z)
\qquad
\text{subject to}
\qquad
\sum_{e\in E}c_eI_e\leq K.
\label{eq:transport-problem}
\end{align}
Let $\mu\geq0$ be the multiplier on the infrastructure budget. The exact KKT conditions are
\begin{align}
\frac{d\widehat U}{dI_e}
&\leq \mu c_e,
&
I_e\left(\mu c_e-\frac{d\widehat U}{dI_e}\right)&=0,
\label{eq:transport-kkt}\\
\mu\left(K-\sum_ec_eI_e\right)&=0.
\label{eq:budget-kkt}
\end{align}
Using \eqref{eq:equilibrium-derivative},
\begin{align}
\frac{d\widehat U}{dI_e}
=
\bar U_{I_e}
-
\bar U_xF_x^{-1}F_{I_e}.
\label{eq:total-investment-gradient}
\end{align}
In the present model, investment affects utility through equilibrium travel costs, so the second term is the central object.

\subsection{Adjoint representation}

Directly computing $F_x^{-1}$ separately for every edge is unnecessary. Let the adjoint vector $\lambda$ solve
\begin{align}
F_x^{\prime}\lambda
=
\bar U_x^{\prime}.
\label{eq:adjoint-system}
\end{align}
Then
\begin{align}
\frac{d\widehat U}{dI_e}
=
\bar U_{I_e}-\lambda^{\prime}F_{I_e}.
\label{eq:adjoint-gradient}
\end{align}
The same adjoint system generates the shadow value of each edge's market access. Let $\Lambda_e\geq0$ denote the resulting shadow-value-weighted importance of edge $e$ across all OD pairs and recursive walks. With the edge technology \eqref{eq:edge-time}--\eqref{eq:edge-cost}, the interior investment condition can be written
\begin{align}
\mu c_e
=
\epsilon\kappa\beta
\delta_e^1n_e^{\rho}\Lambda_e
I_e^{-(\beta+1)}.
\label{eq:investment-foc}
\end{align}
The term $\Lambda_e$ incorporates all equilibrium effects transmitted through OD choices, wages, rents, floor-space supply, amenities, and recursive-path substitution. It must be evaluated at the same equilibrium as $n_e$.

Solving \eqref{eq:investment-foc} gives
\begin{align}
I_e
=
\left(
\frac{\epsilon\kappa\beta}{\mu}
\right)^{1/(\beta+1)}
\left(
\frac{\delta_e^1n_e^{\rho}\Lambda_e}{c_e}
\right)^{1/(\beta+1)}.
\label{eq:investment-level}
\end{align}
Define
\begin{align}
W_e
\equiv
\frac{\delta_e^1n_e^{\rho}\Lambda_e}{c_e}.
\label{eq:edge-weight}
\end{align}
If the budget binds, the optimal allocation satisfies
\begin{align}
\boxed{
I_e^*
=
K
\frac{W_e^{1/(\beta+1)}}
{\displaystyle\sum_{e'\in E}c_{e'}W_{e'}^{1/(\beta+1)}}.
}
\label{eq:investment-rule}
\end{align}
For any two interior edges,
\begin{align}
\frac{I_e^*}{I_{e'}^*}
=
\left(
\frac{W_e}{W_{e'}}
\right)^{1/(\beta+1)}.
\label{eq:investment-ratio}
\end{align}
Because $n_e$ and $\Lambda_e$ depend on $\mathbf I$, \eqref{eq:investment-rule} is a fixed-point characterization of optimal investment, not a one-shot closed form.

\section{Equilibrium with optimal transportation investment}

\begin{definition}[Transport equilibrium under exogenous zoning]
Given zoning $\mathbf z$, a transport equilibrium is a collection
\[
\left(
\mathbf I^*,
\boldsymbol\pi^*,
\mathbf n^*,
\mathbf q^{\R*},
\mathbf q^{\F*},
\mathbf w^*,
\mathbf H^{\R*},
\mathbf H^{\F*}
\right)
\]
such that:
\begin{enumerate}
\item $(\boldsymbol\pi^*,\mathbf n^*)$ solves the spatial fixed point \eqref{eq:spatial-fixed-point} given $(\mathbf I^*,\mathbf z)$;
\item residential and commercial floor-space quantities and prices satisfy \eqref{eq:commercial-quantity-solution}--\eqref{eq:residential-price-solution};
\item firms satisfy \eqref{eq:wage}--\eqref{eq:commercial-rent};
\item household OD and recursive-walk choices satisfy \eqref{eq:route-share}--\eqref{eq:od-share};
\item traffic and travel costs satisfy \eqref{eq:edge-time}--\eqref{eq:edge-traffic};
\item $\mathbf I^*$ solves \eqref{eq:transport-problem}, equivalently satisfying \eqref{eq:transport-kkt}--\eqref{eq:budget-kkt}, and interior edges satisfy \eqref{eq:investment-rule}.
\end{enumerate}
\end{definition}

\section{Solution algorithm}

A nested fixed-point algorithm solves the model:
\begin{enumerate}
\item Fix exogenous zoning $\mathbf z$ and choose an initial feasible investment vector $\mathbf I^{(0)}$.
\item Given $\mathbf I^{(m)}$, solve \eqref{eq:spatial-fixed-point} for $x^{(m)}=(\boldsymbol\pi^{(m)},\mathbf n^{(m)})$. At each inner iteration, use the analytical commercial and residential solutions \eqref{eq:commercial-quantity-solution}--\eqref{eq:residential-price-solution}.
\item Evaluate welfare and solve the adjoint system \eqref{eq:adjoint-system}. Recover $\Lambda_e^{(m)}$ and $W_e^{(m)}$ for every edge.
\item Update investment using
\[
\widetilde I_e^{(m+1)}
=
K
\frac{(W_e^{(m)})^{1/(\beta+1)}}
{\sum_{e'}c_{e'}(W_{e'}^{(m)})^{1/(\beta+1)}}.
\]
Use damping if necessary:
\[
I_e^{(m+1)}
=(1-\omega_I)I_e^{(m)}+\omega_I\widetilde I_e^{(m+1)},
\qquad \omega_I\in(0,1].
\]
\item Repeat steps 2--4 until investment, OD shares, traffic, wages, rents, and floor-space quantities converge.
\end{enumerate}

\section{What exogenous zoning does in this model}

Zoning affects the transportation solution only through the conditional spatial equilibrium. A tighter residential FAR cap lowers $\bar H_i^{\R}$. When binding, it fixes local housing quantity and causes demand changes to appear primarily in $q_i^{\R}$. This changes residential sorting, OD flows, edge traffic, and the edge weights $W_e$. A tighter commercial cap similarly changes local employment capacity, wages, workplace choice, and traffic.

With a fixed infrastructure budget $K$, exogenous zoning changes the allocation of investment across edges. It cannot change aggregate investment expenditure. A model of aggregate overinvestment or underinvestment requires replacing the fixed budget with an endogenous investment-cost problem such as
\[
\max_{\mathbf I\geq0}
\widehat U(\mathbf I;\mathbf z)-\sum_eC_e(I_e).
\]
Under the current fixed-budget formulation, the correct outcome is spatial misallocation of a given transportation budget across links.

\printbibliography

\end{document}
